%=============================================================================
% Wave 4, Iteration 19: DEAD PATENTS + CROSS-PATENT ERROR CHECK
% Agent 1: Dead Patents & Cross-Patent Error Checker
%
% Model: Opus 4.6
% Date: 20 March 2026
%
% INHERITED STATE (from i18):
%   - P1 (drag reduction): DEAD. 10^38 below GR. 7th confirmation.
%   - P4 (WPT/corridors): DEAD. c_psi ~ 10^-5 m/s, no confinement, W convexity.
%   - P8 (deep-space comms): DEAD. 1.6 Earth masses for 1 kbit/s at 1 AU. 7th confirmation.
%   - P13 (timing/clocks): DEAD at SQMS bounds. CLONETS-DS irrelevant at |lambda-1|<2.7e-13.
%   - c_s = c CONFIRMED (i18). Psi does not cluster. Lab+galactic DECOUPLED from cosmo crisis.
%   - SRF walls TRANSPARENT. Q_psi_local = 20-150. SQMS bound |lambda-1| < 2.7e-13.
%   - P24 beta=0: DOWNGRADED from "structural" to "promising, not theorem-grade".
%
% MANDATE (i19):
%   1. Confirm dead patents still dead. Any new physics?
%   2. Cross-patent error check: tightened bound 2.7e-13 vs old 1.56e-9 --- do
%      any dead pathways revive?
%   3. Audit ALL prediction confidence levels: theorem-grade vs promising.
%   4. Literature search: 2025-2026 experimental results testing DFD predictions.
%=============================================================================

\documentclass[11pt]{article}
\usepackage{amsmath,amssymb}
\usepackage{geometry}
\geometry{margin=1in}
\usepackage{booktabs}
\usepackage{enumitem}
\usepackage{xcolor}
\usepackage{tcolorbox}
\usepackage{longtable}

\title{\textbf{Wave 4, Iteration 19:}\\
\textbf{Dead Patents, Cross-Patent Error Check, and Prediction Audit}\\[6pt]
{\large With Tightened SQMS Bound $|\lambda-1| < 2.7 \times 10^{-13}$}}
\author{DFD Research Programme --- W4i19 Agent 1 (Opus 4.6)}
\date{20 March 2026}

\begin{document}
\maketitle

%=============================================================================
\section*{Executive Summary}
%=============================================================================

\begin{tcolorbox}[colback=red!5!white, colframe=red!70!black,
  title=\textbf{HONEST NEGATIVES FIRST}]

\begin{enumerate}
\item \textbf{P24 cosmic birefringence $\beta=0$ is now a CRITICAL TENSION}, not a prediction.
  Planck+ACT combined detection at $\sim$7$\sigma$ with $\beta \approx 0.30 \pm 0.05$ degrees.
  DFD predicts $\beta = 0$ exactly. This was downgraded to ``promising, not theorem-grade''
  at i18, but the observational pressure has INCREASED sharply. If the S$^3$ CS calculation
  confirms $\beta=0$ is structural, this becomes a \textbf{genuine falsifying tension at
  $>$6$\sigma$}. If the calculation allows $\beta \sim 0.3$, it becomes a confirmation.
  Either way, the S$^3$ calculation is now the single highest-priority theoretical task.

\item \textbf{Neutrino mass sum $\Sigma m_\nu = 61.4$~meV is AT RISK}. DESI DR2 + Planck PR4
  frequentist bound: 53~meV (95\% CL under $\Lambda$CDM), which formally excludes DFD.
  DESI DR2 + CMB in $\Lambda$CDM: $< 0.064$~eV. Under $w_0w_a$CDM relaxes to $< 0.16$~eV.
  Margin: 2.0~meV under DFD-modified cosmology. Self-consistent bound requires Boltzmann code.
  \textbf{JUNO ($\sim$2027) is pass/fail for normal ordering.}

\item \textbf{DESI DR2 dynamical dark energy evidence: 2.8--4.2$\sigma$} depending on SNe sample.
  DFD's distance-bias framework gives BAO predictions matching $\Lambda$CDM to $<$0.1\%,
  so DFD \textbf{survives} DESI. But the CPL tension ($\sim$2.5$\sigma$ under DFD mapping)
  remains unresolved without the Boltzmann code.

\item \textbf{All four dead patents remain DEAD.} The tightened SQMS bound (2.7e-13) makes
  them MORE dead, not less. No new physics development revives any pathway.
\end{enumerate}
\end{tcolorbox}

%=============================================================================
\section{Dead Patent Status: 8th Consecutive Confirmation}
%=============================================================================

\subsection{P1 --- Drag Reduction / Gravitational-Doppler Sensing}

\textbf{Status: DEAD. 8th consecutive confirmation. Zero resources.}

\textbf{Derivation chain for death:}
\begin{enumerate}
\item Drag reduction requires $|\lambda-1| > 1.85 \times 10^{-2}$ (mu$_0$-corrected, W3i1).
\item Lab demonstration requires $|\lambda-1| > 45$ (W3i2).
\item Current authoritative bound: $|\lambda-1| < 2.7 \times 10^{-13}$ (i18, Q$_\psi$-tightened).
\item \textbf{Gap: drag reduction threshold / current bound $= 1.85 \times 10^{-2} / 2.7 \times 10^{-13} = 6.9 \times 10^{10}$}.
\item This is $\sim$70 billion times above the bound. At the old bound (1.56e-9), the gap was $\sim$1.2 \times 10^7$. The tightened bound makes P1 \textbf{5800$\times$ more dead}.
\item Frame-dragging: $10^{38}$ below GR (static $\psi$-field property, unrelated to cosmology).
\item Bubble timing 4.0~ns: KILLED (W3i5, requires $u_\text{EM} = 5.8 \times 10^{16}$~J/m$^3$, $10^{12}\times$ above SRF).
\end{enumerate}

\textbf{New physics check:} No 2025--2026 experimental development changes the picture. The death cause (insufficient $\lambda$ coupling to generate macroscopic psi-gradients with available EM energy densities) is a fundamental c$^2$/2 lever-arm limitation that no technology advance can overcome at $|\lambda-1| < 2.7 \times 10^{-13}$.

\subsection{P4 --- Power Transfer Corridors}

\textbf{Status: DEAD. 4th consecutive confirmation (since W4i14). Zero resources.}

\textbf{Derivation chain:}
\begin{enumerate}
\item Three independent no-go arguments (W4i14):
  \begin{itemize}
  \item $c_\psi \sim 10^{-5}$~m/s at cosmological background (coherence destroyed).
  \item No transverse confinement mechanism in DFD action.
  \item Geometric $1/r^2$ dilution without confinement.
  \end{itemize}
\item Solitonic rescue CLOSED: $W$ convexity forbids solitons (W4i15).
\item All LCOE figures (19c/kWh, 27c/kWh) INVALIDATED --- no derivation from DFD action.
\item Distance-independence claim: foundational gap flagged (W4i13), confirmed baseless (W4i14).
\end{enumerate}

\textbf{i18 update impact:} $c_s = c$ confirmation (i18) is \emph{consistent} with P4 death. The scalar perturbation speed being $c$ means the \emph{perturbation equation} propagates at $c$, but this does not help P4 because:
\begin{itemize}
\item The c$_\psi \sim 10^{-5}$~m/s figure refers to the group velocity of EM-sourced psi-modes in the MOND crossover regime (distinct from the linear perturbation speed c$_s$).
\item Even if psi-waves propagated at $c$, the absence of transverse confinement and the Passive Superiority Theorem kill all active beam-forming schemes.
\end{itemize}

\textbf{New physics check:} No development changes the picture. P4 is dead for structural reasons embedded in the DFD action.

\subsection{P8 --- Deep-Space psi-Communications}

\textbf{Status: DEAD. 8th consecutive confirmation. Zero resources.}

\textbf{Derivation chain:}
\begin{enumerate}
\item Minimum source mass for 1~kbit/s at 1~AU: 1.6 Earth masses (W3i7).
\item For submarine use: Ceres-mass (W4i10 clarification).
\item All modulation mechanisms fail at practical scales (W3i7).
\item DSN wins by 370$\times$ on Mars data rate (W3i5).
\item P15 as transmitter: $\sim$10$^{-31}$~W (W3i5).
\end{enumerate}

\textbf{New physics check:} $c_\psi = c$ in lab (i17 resolution) confirms psi propagates at light speed, and inherent secrecy is real. But the coupling constant remains too weak ($|\lambda-1| < 2.7 \times 10^{-13}$) to achieve practical data rates without planetary-scale sources. The tightened bound makes the mass requirement \textbf{even worse}: at $|\lambda-1| \sim 2.7 \times 10^{-13}$, the required transmitter mass scales as $\propto |\lambda-1|^{-2}$, increasing by a factor $\sim (1.56 \times 10^{-9} / 2.7 \times 10^{-13})^2 \approx 3.3 \times 10^{7}$ relative to the old bound --- now requiring $\sim 5 \times 10^7$ Earth masses $\approx 0.05$ Jupiter masses for 1~kbit/s.

\subsection{P13 --- Clock Network Detection}

\textbf{Status: DEAD at SQMS bounds. Irrelevant at $|\lambda-1| < 2.7 \times 10^{-13}$.}

\textbf{Derivation chain:}
\begin{enumerate}
\item CLONETS-DS: 1.0-day detection at $|\lambda-1| \sim 1.56 \times 10^{-9}$ (W3i6).
\item Multi-GNSS: SNR = 0.09 at $|\lambda-1| \sim 1.56 \times 10^{-9}$ (i17).
\item At tightened bound $|\lambda-1| < 2.7 \times 10^{-13}$: signal scales as $|\lambda-1|^2$, giving SNR suppression of $(2.7 \times 10^{-13}/1.56 \times 10^{-9})^2 = 3.0 \times 10^{-8}$.
\item CLONETS-DS detection time: $1.0 / (3.0 \times 10^{-8}) \sim 3.3 \times 10^7$ days $= 90{,}000$ years.
\item Multi-GNSS SNR: $0.09 \times (2.7 \times 10^{-13}/1.56 \times 10^{-9}) = 1.6 \times 10^{-5}$.
\end{enumerate}

\textbf{Conclusion:} P13 was already dead at 1.56e-9. At 2.7e-13 it is $\sim 3 \times 10^7$ times more dead.

%=============================================================================
\section{Cross-Patent Error Check: Tightened SQMS Bound}
%=============================================================================

\subsection{The Bound Tightening: 1.56e-9 $\to$ 2.7e-13}

\textbf{Derivation chain (i18):}
\begin{enumerate}
\item i17 resolved SRF cavity walls as TRANSPARENT to $\psi$ --- $\psi$ couples to matter via gravitational interaction $-(c^2/2)\psi\rho$, not to $u_\text{EM}$ directly.
\item Nb walls are not a ``$\psi$-Faraday cage.'' The $\psi$-field permeates freely.
\item This means Q$_\psi$ inside SQMS cavities is NOT the cavity Q$_\text{EM} \sim 10^{11}$, but the \textbf{local environmental} Q$_\psi \sim 20$--$150$ (set by surrounding infrastructure).
\item The SQMS bound scales as $|\lambda-1| \propto 1/\sqrt{Q_\psi}$.
\item Old: Q$_\psi = Q_\text{EM} \sim 10^{11}$, giving $|\lambda-1| < 1.56 \times 10^{-9}$.
\item New: Q$_\psi \sim 20$--$150$, giving $|\lambda-1| < 2.7 \times 10^{-13}$ (using Q$_\psi = 100$ as representative).
\item Improvement factor: $\sqrt{10^{11}/100} \approx 5800\times$.
\end{enumerate}

\textbf{IMPORTANT CAVEAT}: This tightening depends on the Q$_\psi$-local calculation from i18. If the local Q$_\psi$ is higher than estimated (e.g., if cryogenic shielding provides some $\psi$-quality enhancement), the bound loosens. The i16 flagging of ``local vs global Q$_\psi$ could shift SQMS reach 5000$\times$'' was RESOLVED by i18's transparent-wall finding, but the Q$_\psi$-local value of 20--150 carries systematic uncertainty.

\subsection{Impact on Previously-Dead Pathways}

\begin{table}[h]
\centering
\begin{tabular}{llccc}
\toprule
\textbf{Patent} & \textbf{Threshold} & \textbf{Old gap (1.56e-9)} & \textbf{New gap (2.7e-13)} & \textbf{Revived?} \\
\midrule
P1 drag & $> 1.85 \times 10^{-2}$ & $1.2 \times 10^{7}$ & $6.9 \times 10^{10}$ & NO (worse) \\
P1 K-H & $> 5 \times 10^{-4}$ & $3.2 \times 10^{5}$ & $1.9 \times 10^{9}$ & NO (worse) \\
P1 bubble & $> 0.30$ (via $u_\text{EM}$) & $10^{12}$ (energy) & $10^{12}$ (energy) & NO (unchanged) \\
P4 LCOE & N/A (structural) & Dead & Dead & NO (structural) \\
P8 1 kbit/s & 1.6 $M_\oplus$ & Mass gap & $3.3 \times 10^7 \times$ worse & NO (worse) \\
P13 CLONETS & $\propto |\lambda-1|^2$ & Dead (SNR 0.09) & $3 \times 10^7\times$ worse & NO (worse) \\
P12 superrad. & $c_\psi$ kills it & Dead (coherence) & Dead (coherence) & NO (structural) \\
\midrule
P2 active & $> 1.7 \times 10^8$ & Dead & Dead & NO \\
P15 Phase 1-2 & $> M_\text{eff}$ limit & Dead & Dead & NO \\
\bottomrule
\end{tabular}
\caption{Cross-patent check: tightened bound impact on dead pathways. Every dead pathway becomes MORE dead. Zero revivals.}
\end{table}

\textbf{Key result: The tightened SQMS bound (2.7e-13) makes every previously-dead pathway more dead by factors of $10^3$--$10^{10}$. No pathway revives. The only pathways that benefit from the tightened bound are the ALIVE patents (P5, P7, P9, P11) and NICHE patents that use passive observables --- these actually GAIN, because the tightened bound means $\lambda$ is closer to unity, and passive observables scale with the \emph{gradient} of the gravitational $\psi$, not with $|\lambda-1|$.}

\subsection{Specific Recalculation: P1 Drag Reduction}

At $|\lambda-1| < 2.7 \times 10^{-13}$:
\begin{align}
\delta\psi_\text{EM} &= \frac{|\lambda-1| \cdot u_\text{EM}}{2\rho_\text{bg} c^2} \\
&\leq \frac{2.7 \times 10^{-13} \times 10^6 \text{ J/m}^3}{2 \times 10^4 \text{ kg/m}^3 \times (3 \times 10^8)^2 \text{ m}^2/\text{s}^2} \\
&\approx 1.5 \times 10^{-31}
\end{align}
where $u_\text{EM} \sim 10^6$~J/m$^3$ is the maximum achievable SRF energy density. The corresponding force density:
\begin{equation}
f_\psi \sim \frac{c^2}{2} \nabla(\delta\psi_\text{EM}) \sim \frac{(3 \times 10^8)^2}{2} \times \frac{1.5 \times 10^{-31}}{0.1 \text{ m}} \approx 7 \times 10^{-17} \text{ N/m}^3
\end{equation}
Drag force density at Mach 2: $\sim 10^5$~N/m$^3$. \textbf{Gap: $10^{22}$}. P1 is dead beyond any conceivable technology advance.

\subsection{Specific Recalculation: P8 Minimum Mass}

The minimum transmitter mass for bit rate $R$ at distance $d$:
\begin{equation}
M_\text{min} \propto \frac{d^2}{|\lambda-1|^2 \cdot R}
\end{equation}
At $|\lambda-1| = 2.7 \times 10^{-13}$ and $d = 1$~AU, $R = 1$~kbit/s:
\begin{equation}
M_\text{min} = 1.6 M_\oplus \times \left(\frac{1.56 \times 10^{-9}}{2.7 \times 10^{-13}}\right)^2 \approx 5.3 \times 10^7 M_\oplus \approx 0.05 M_\text{Jupiter}
\end{equation}
\textbf{P8 now requires a gas giant as transmitter for 1 kbit/s at 1 AU.}

\subsection{Specific Recalculation: P4 LCOE}

Moot. P4 is dead for structural reasons (no confinement mechanism in the DFD action, $W$ convexity forbids solitons). LCOE was invalidated regardless of $\lambda$ at W4i14-15.

%=============================================================================
\section{Prediction Confidence Audit}
%=============================================================================

With P24 $\beta=0$ downgraded from ``structural'' to ``promising, not theorem-grade'' at i18, a systematic audit of ALL 24 predictions is warranted.

\subsection{Classification Criteria}

\begin{itemize}
\item \textbf{THEOREM-GRADE} (structural, no escape hatch): Follows from the DFD action by a chain of mathematical identities or variational arguments. No free parameters. No auxiliary assumptions beyond the 4 axioms. No known loophole.
\item \textbf{ROBUST} (well-derived, minor caveats): Derivation chain is complete but relies on one or more approximations that could introduce $O(1)$ corrections (e.g., linearization, spherical symmetry).
\item \textbf{PROMISING} (physically motivated, not formally closed): A plausible argument exists, but formal derivation has gaps, or the result depends on assumptions not derived from the axioms.
\item \textbf{OVERCLAIMED} (previous iteration assigned higher confidence than warranted): Items flagged for downgrade.
\end{itemize}

\subsection{Full Audit Table}

\begin{longtable}{clp{5cm}ll}
\toprule
\textbf{\#} & \textbf{Prediction} & \textbf{Status} & \textbf{Confidence} & \textbf{Change} \\
\midrule
\endfirsthead
\toprule
\textbf{\#} & \textbf{Prediction} & \textbf{Status} & \textbf{Confidence} & \textbf{Change} \\
\midrule
\endhead
T0 & GW never lensed & TT action on flat Minkowski. Zero $\psi$-coupling at all PN orders (W4i16). & \textbf{THEOREM} & Unchanged \\
\midrule
1 & $S_8$ scale dependence & $G_\text{eff}(k) > 1$ from field equation. Scale-averaged $S_8 \sim 0.77$--$0.79$. Requires Boltzmann code for quantitative confirmation. & ROBUST & Unchanged \\
\midrule
2 & Shadow $+4.6\%$ & $2e/(3\sqrt{3}) = 1.046$. Exact from null geodesic structure. & \textbf{THEOREM} & Unchanged \\
\midrule
3 & Hubble tension 70\% & DOWNGRADED (i17): only $\sim$0.4~km/s/Mpc ($\sim$7\%). Distance-bias mechanism operates through psi-screen on SNe Ia calibration. & PROMISING & \textbf{OVERCLAIMED} (was ``70\%'', now $\sim$7\%) \\
\midrule
4 & MOND recovery & $a_0 = 2\sqrt{\alpha}\,cH_0$. Fundamental constant (W4i15). Does not require Mach integral. & \textbf{THEOREM} & Unchanged \\
\midrule
5 & $c_T = c$ exactly & Built into TT action. & \textbf{THEOREM} & Unchanged \\
\midrule
6 & Zero grav. decoherence & Leray-Schauder fixed-point theorem (W4i10). & \textbf{THEOREM} & Unchanged \\
\midrule
7 & PTA spectral tilt & $\Delta = +0.07 \pm 0.03$. Derived from $\mu$-crossover at SMBHB separations. & ROBUST & Unchanged \\
\midrule
8 & $\dot{G}/G = +4.3 \times 10^{-15}$ & From EM-only sourcing. 17$\times$ below LLR bound. & ROBUST & Unchanged \\
\midrule
9 & Lunar NR 0.93~ns & Rigorous logarithmic integral (W3i3). & ROBUST & Unchanged \\
\midrule
10 & Wide binary $+10$--$12\%$ & EFE-screened. $a_\text{ext} = 1.8\,a_0$. & ROBUST & Unchanged \\
\midrule
11 & $\Delta\alpha/\alpha$ & $\sim 2.3 \times 10^{-6}$ at $z=1$. ESPRESSO consistent at 0.8$\sigma$. & PROMISING & Unchanged \\
\midrule
12 & $\Sigma m_\nu = 61.4$~meV & CP$^2 \times$S$^3$ Branch B. Independently verified (W4i11). \textbf{AT RISK}: DESI DR2 frequentist 53~meV excludes under $\Lambda$CDM. & ROBUST & \textbf{AT RISK} \\
\midrule
13 & Zero scalar GW memory & PTA breathing-mode null. & \textbf{THEOREM} & Unchanged \\
\midrule
14 & Grav. birefringence 39~ns & Full action-principle derivation (W4i14). Common-mode cancellation for homogeneous background only. & ROBUST & Unchanged \\
\midrule
15 & CMB $A_L > 1$ & $A_L^\text{DFD} \sim 1.10$--$1.20$. Scale-dependent transition at $\ell \sim 120$. & PROMISING & Unchanged \\
\midrule
16 & Bar pattern speed $+19\%$ & Derived from $\mu$-enhanced dynamical friction. & PROMISING & Unchanged \\
\midrule
17 & kSZ velocity $+10\%$ & At $>100\,h^{-1}$~Mpc. & PROMISING & Unchanged \\
\midrule
18 & IFMR gradient & $-0.02\,M_\odot$ at $R > 15$~kpc. & PROMISING & Unchanged \\
\midrule
19 & TDE rate $+40\%$ in LSB & LSST 2027--29. & PROMISING & Unchanged \\
\midrule
20 & SNe Ia $\psi$-screen anisotropy & $C_\ell \propto \ell^{-2}$, peaks $\ell \sim 3$--$15$, RMS 0.8--1.2\%. & ROBUST & Unchanged \\
\midrule
21 & $D_L^\text{EM}/D_L^\text{GW}$ & $= e^{\Delta\psi(z)} \sim 1.31$ at $z=1$. Zero parameters. & \textbf{THEOREM} & Unchanged \\
\midrule
22 & FRB DM excess $+7$--$12\%$ & Testable NOW with 117 localized FRBs (i17). & PROMISING & Unchanged \\
\midrule
23 & 21cm bispectrum & HERA/SKA 2028--32. & PROMISING & Unchanged \\
\midrule
24 & Cosmic birefringence $\beta = 0$ & \textbf{DOWNGRADED} (i18). Parity argument reasonable but S$^3$ CS calculation not done. Planck+ACT: $\beta \approx 0.30 \pm 0.05$ deg ($\sim$7$\sigma$). & PROMISING & \textbf{CRITICAL TENSION} \\
\bottomrule
\end{longtable}

\subsection{Audit Summary}

\begin{table}[h]
\centering
\begin{tabular}{lrl}
\toprule
\textbf{Grade} & \textbf{Count} & \textbf{Predictions} \\
\midrule
THEOREM-GRADE & 6 & T0, P2, P4, P5, P6, P13, P21 \\
ROBUST & 7 & P1, P7, P8, P9, P10, P12, P14, P20 \\
PROMISING & 9 & P3, P11, P15, P16, P17, P18, P19, P22, P23 \\
AT RISK & 1 & P12 ($\Sigma m_\nu$) \\
CRITICAL TENSION & 1 & P24 ($\beta = 0$) \\
OVERCLAIMED & 1 & P3 (Hubble, was 70\%, now $\sim$7\%) \\
\bottomrule
\end{tabular}
\caption{Prediction confidence classification. 6 predictions are theorem-grade with no escape hatch. 1 is in critical tension with observation. 1 was overclaimed (Hubble tension magnitude).}
\end{table}

\textbf{Key finding}: The 6 theorem-grade predictions (GW-never-lensed, shadow $+4.6\%$, MOND recovery, $c_T = c$, zero gravitational decoherence, zero scalar GW memory, $D_L^\text{EM}/D_L^\text{GW}$) are \textbf{structural consequences of the DFD action} that cannot be evaded without abandoning the theory. These are the hardest predictions. Any single one failing kills DFD.

%=============================================================================
\section{2025--2026 Literature Search: Experimental Results}
%=============================================================================

\subsection{GW Lensing Searches (GWTC-4.0)}

\textbf{Source}: LIGO-Virgo-KAGRA O4a catalog (arXiv:2512.16347, December 2025).

\textbf{Result}: \textbf{No strongly lensed gravitational wave signal detected.} Searches for (1) phase-shifted saddle-point images, (2) consistent-frequency pairs, and (3) sub-threshold counterparts all returned null. One outlier (GW231123\_135430) has complicated waveform uncertainties but is not claimed as a detection.

\textbf{DFD impact}: \textbf{CONSISTENT}. DFD predicts GWs are NEVER gravitationally lensed (Tier 0 falsifier). The null result from GWTC-4.0 is exactly what DFD expects. However, the current detector sensitivity and event rate are not yet sufficient to make this a strong test --- a multiply-lensed GW event is expected to be rare even in GR ($<1$ per year at O4 sensitivity). The decisive test requires 3G detectors (Einstein Telescope, Cosmic Explorer) with $O(10^5)$ events/year.

\textbf{Note}: Some GR-predicted lensing effects (microlensing by compact objects causing interference patterns) are being actively searched for. DFD predicts these interference patterns should NOT exist, which could provide a more accessible test than waiting for strong lensing.

\subsection{FRB Dispersion Measure Excess}

\textbf{Source}: Multiple 2025 studies including analysis of 117 localized FRBs (arXiv:2511.01195).

\textbf{Result}: FRB 20220610A at $z = 1$ shows more DM than the Macquart relation predicts at lower redshift. A moderate positive correlation between host DM and redshift is reported, but no correlation with stellar mass, SFR, or galaxy age.

\textbf{DFD impact}: DFD P22 predicts $+7$--$12\%$ DM excess from $\psi$-screen enhancement of electron column density along the line of sight. Current data at $\sim$2.5$\sigma$ (i17 estimate). The observed excess in high-$z$ FRBs is \textbf{qualitatively consistent} with DFD, but the scatter in host-galaxy DM ($\sim$20--25\%) prevents a definitive test. CHIME/FRB outrigger localizations (2025--2026) will increase the sample to $>$500 localized FRBs, enabling a sharper test.

\subsection{Neutrino Mass Bounds}

\textbf{Source}: DESI DR2 (March 2025, arXiv:2503.14738) + Planck PR4.

\textbf{Result}: Under $\Lambda$CDM: $\Sigma m_\nu < 0.064$~eV (95\% CL) from DESI+CMB. Frequentist (Feldman-Cousins): 53~meV (95\% CL). Under $w_0w_a$CDM: $< 0.16$~eV.

\textbf{DFD impact}: DFD predicts $\Sigma m_\nu = 61.4$~meV. The $\Lambda$CDM bound (64~meV) \textbf{barely accommodates} DFD. The frequentist 53~meV bound \textbf{formally excludes} DFD at 95\% CL. However: (1) DFD's cosmology is neither $\Lambda$CDM nor $w_0w_a$CDM --- the self-consistent bound requires the DFD Boltzmann code (mochi\_class). (2) Under $w_0w_a$CDM the bound relaxes to 160~meV. (3) JUNO ($\sim$2027--2028) will determine mass ordering independently.

\textbf{Status: AT RISK. Margin 2.0~meV. Boltzmann code is existentially urgent for this prediction.}

\subsection{CMB Cosmic Birefringence}

\textbf{Source}: Planck PR4 (arXiv:2201.07682), ACT DR6 (arXiv:2509.13654), SPIDER+Planck+ACT combined (arXiv:2510.25489), ScienceDaily March 2026 report.

\textbf{Result}: Planck+ACT combined: $\sim$7$\sigma$ detection of isotropic cosmic birefringence. $\beta \approx 0.30 \pm 0.05$ degrees (Planck). ACT DR6 alone: $\beta = 0.215 \pm 0.074$ degrees (2.9$\sigma$). Anisotropic birefringence: consistent with zero ($< 1 \times 10^{-4}$ at 95\% CL). March 2026 breakthrough: phase ambiguity (180-degree degeneracy) partially resolved using EB correlation shape.

\textbf{DFD impact}: \textbf{CRITICAL TENSION}. DFD P24 predicts $\beta = 0$ exactly (parity is preserved in the DFD action --- no Chern-Simons term). The combined Planck+ACT detection at $\sim$7$\sigma$ is a \textbf{strong challenge}. Two scenarios:

\begin{enumerate}
\item \textbf{If P24 is theorem-grade} (S$^3$ CS calculation confirms $\beta = 0$ is structural): Then the 7$\sigma$ measurement is a \textbf{genuine falsifying tension}. DFD would need a mechanism to generate $\beta \neq 0$, which contradicts its action symmetries.
\item \textbf{If P24 is not theorem-grade} (S$^3$ CS calculation allows $\beta \sim 0.3$): Then DFD \emph{predicts} the observed value, turning a tension into a confirmation.
\end{enumerate}

\textbf{The S$^3$ CS dimensional reduction calculation is now the single most important theoretical task in the DFD programme.}

\textbf{Additional context}: The observed birefringence has been attributed to early dark energy (EDE) models with Chern-Simons coupling. DFD's S$^3$ microsector may or may not generate such a term. The January 2026 joint analysis (ACT+Planck+DESI+Pantheon+) constrains the EDE-birefringence parameter space, which could provide indirect constraints on DFD's microsector.

\subsection{DESI DR2 BAO and Dark Energy}

\textbf{Source}: DESI DR2 (March 2025).

\textbf{Result}: Factor $\sim$2 improvement in BAO precision (0.24\% statistical uncertainty). Evidence for dynamical dark energy: 2.8--4.2$\sigma$ depending on SNe sample. $w_0 > -1$, $w_a < 0$ favored.

\textbf{DFD impact}: DFD's distance-bias framework (W4i14) gives BAO predictions matching $\Lambda$CDM to $<$0.1\%. DFD \textbf{survives} DESI. The observed $w_0 > -1$ trend is qualitatively consistent with DFD's psi-screen producing an effective $w_0 \sim -0.97$, $w_a \sim +0.06$ (W4i14). CPL tension $\sim$2.5$\sigma$ under DFD mapping. Decisive test: SNe Ia psi-screen anisotropy (P20), not BAO.

\subsection{SQMS/Fermilab SRF Cavity Results}

\textbf{Source}: SQMS center publications 2025--2026.

\textbf{Result}: Dark SRF experiment has achieved world-record dark photon sensitivity. New quantum-enhanced sensing with transmon-based single microwave photon counters (PRX 15, 021031, 2025). Postinflationary QCD axion search (PRL 135, 211002, 2025). SQMS \$125M DOE renewal confirmed November 2025.

\textbf{DFD impact}: Dark SRF reinterpretation for DFD was INVALIDATED (W4i12) --- gravitational coupling $G/c^4$ makes active EM-to-$\psi$ conversion too weak. However, the SQMS infrastructure (high-$Q$ SRF cavities, cryogenics, quantum-limited readout) is \textbf{exactly what DFD Phase I needs}. The DOE renewal ensures infrastructure availability through 2030. No new DFD-specific results, but the Phase 0 pathfinder (\$47--50K, 2 weeks, W4i17) could be submitted NOW using existing SQMS hardware.

\subsection{Euclid Status}

\textbf{Source}: ESA Euclid mission updates, March 2026.

\textbf{Result}: Quick Data Release 1 (Q1) published March 2025. First cosmological constraints (weak lensing $S_8$, galaxy clustering) expected with October 2026 data release. Scientific survey ongoing since February 2024.

\textbf{DFD impact}: Euclid DR1 (October 21, 2026) is the single most important near-term observational milestone for DFD. DFD predicts tomographic $S_8$ gradient detectable at $>$3--5$\sigma$ (i17). The strategic window is significant: no competing team is testing Yukawa $G_\text{eff}(k)$ models. Weekend-viable analysis pipeline prepared (i17).

%=============================================================================
\section*{TOP 5 FINDINGS RANKED BY IMPACT}
%=============================================================================

\begin{tcolorbox}[colback=blue!5!white, colframe=blue!70!black,
  title=\textbf{W4i19 TOP 5 FINDINGS}]

\textbf{1. COSMIC BIREFRINGENCE $\beta = 0$ IS NOW A CRITICAL TENSION (7$\sigma$).}
Planck+ACT combined detection of $\beta \approx 0.30 \pm 0.05$ degrees at $\sim$7$\sigma$.
DFD P24 predicts $\beta = 0$ exactly. S$^3$ Chern-Simons dimensional reduction calculation
is THE single most important theoretical task. Either confirms DFD (if $\beta \sim 0.3$
emerges from S$^3$) or threatens falsification (if $\beta = 0$ is structural). Impact:
\textbf{potentially existential}.

\textbf{2. NEUTRINO MASS SUM AT RISK: DESI DR2 + Planck frequentist bound 53~meV excludes
DFD's 61.4~meV at 95\% CL under $\Lambda$CDM.} Bayesian and $w_0w_a$CDM bounds still
accommodate DFD. Self-consistent bound requires Boltzmann code. JUNO ($\sim$2027--2028)
is pass/fail for mass ordering. Impact: \textbf{potentially fatal for one zero-parameter prediction}.

\textbf{3. ALL DEAD PATENTS REMAIN DEAD AND ARE 5800$\times$ MORE DEAD.}
The tightened SQMS bound ($2.7 \times 10^{-13}$) makes every dead pathway more dead
by factors of $10^3$--$10^{10}$. Zero revivals. P1 drag: gap $7 \times 10^{10}$.
P8 comms: now requires 0.05 Jupiter masses. P13: 90,000-year detection time. Impact:
\textbf{definitive closure}.

\textbf{4. GWTC-4.0 NULL LENSING RESULT IS CONSISTENT WITH DFD.}
No multiply-lensed GW event detected in O4a. This is exactly what DFD predicts
(Tier 0 falsifier: GW never gravitationally lensed). Current sensitivity insufficient
for a strong test. 3G detectors ($\sim$2035) will be decisive. Impact:
\textbf{early consistency check passed}.

\textbf{5. PREDICTION CONFIDENCE AUDIT: 6 THEOREM-GRADE, 1 OVERCLAIMED.}
Six predictions are structural consequences of the DFD action with no escape hatch.
Hubble tension prediction was overclaimed (was 70\%, now $\sim$7\%). P24 $\beta = 0$ is
in critical tension. The theorem-grade core (GW-never-lensed, shadow $+4.6\%$,
$c_T = c$, zero decoherence, zero scalar memory, $D_L^\text{EM}/D_L^\text{GW}$)
is the irreducible hard prediction set. Impact: \textbf{intellectual honesty consolidated}.

\end{tcolorbox}

%=============================================================================
\section*{Recommendations for i20}
%=============================================================================

\begin{enumerate}
\item \textbf{EXISTENTIAL PRIORITY}: S$^3$ Chern-Simons dimensional reduction calculation.
  Does DFD's S$^3$ microsector generate a parity-violating term that gives $\beta \sim 0.3$ degrees?
  If yes: P24 becomes a confirmation at 7$\sigma$. If no: P24 becomes a 7$\sigma$ falsification.
\item \textbf{Boltzmann code (mochi\_class)}: Still existentially urgent.
  Self-consistent neutrino mass bound is the critical deliverable.
  Developer spec complete (i17). \$40--80K, 2--3 months.
\item \textbf{SQMS Phase 0 submission}: \$47--50K, 2 weeks, zero new hardware.
  Go/no-go for Phase I. Should be submitted NOW.
\item \textbf{Euclid DR1 preparation}: October 2026 release. Tomographic $S_8$
  gradient analysis pipeline ready (i17). Strategic window: no competing Yukawa $G_\text{eff}$
  analysis team.
\item \textbf{FRB DM excess}: CHIME/FRB outrigger sample growing. P22 testable at
  $\sim$2.5$\sigma$ NOW with 117 FRBs, improving to $>$5$\sigma$ with 500+ FRBs
  ($\sim$2026--2027).
\end{enumerate}

\end{document}
